Graf odchylek sektorových mezd od celostátního mediánu (pro CZ vs.~AT a~DK) kvantifikuje míru polarizace trhu práce: v~ČR jsou odchylky výrazně větší --- ICT a~finance vykazují prémii +\SIrange{60}{80}{\percent} nad mediánem, zatímco obchod a~pohostinství zaostávají o~\SIrange{30}{40}{\percent}.
V~zemích s~erga omnes kolektivním vyjednáváním (AT, DK) je inter-sektorová disperze výrazně nižší: sektorové smlouvy s~celoodvětvovou platností zakotvují spodní sektorový standard, který redukuje rozsah, do jakého může spodní sektor klesat pod celostátní průměr.
Vysoká inter-sektorová disperze v~ČR má reálné sociální důsledky: pracovníci v~obchodu a~pohostinství --- skupiny s~vysokým podílem žen a~nízkokvalifikovaných --- jsou mzdově vytlačeni do in-work poverty territory (srov.~obr.~\ref{fig:arope_groups}).
Zmenšení sektorových odchylek směrem k~AT nebo DK modelu je proto přímo cílem doporučené reformy: erga omnes s~reálnými minimálními tarify v~každém \ac{NACE} odvětví by stlačilo dolní odchylku bez nutnosti administrativního zasahování do horního pásma.

Graf odchylek sektorových mezd od celostátního mediánu (pro CZ vs.~AT a~DK) kvantifikuje míru polarizace trhu práce: v~ČR jsou odchylky výrazně větší --- ICT a~finance vykazují prémii +\SIrange{60}{80}{\percent} nad mediánem, zatímco obchod a~pohostinství zaostávají o~\SIrange{30}{40}{\percent}.
V~zemích s~erga omnes kolektivním vyjednáváním (AT, DK) je inter-sektorová disperze výrazně nižší: sektorové smlouvy s~celoodvětvovou platností zakotvují spodní sektorový standard, který redukuje rozsah, do jakého může spodní sektor klesat pod celostátní průměr.
Vysoká inter-sektorová disperze v~ČR má reálné sociální důsledky: pracovníci v~obchodu a~pohostinství --- skupiny s~vysokým podílem žen a~nízkokvalifikovaných --- jsou mzdově vytlačeni do in-work poverty territory (srov.~obr.~\ref{fig:arope_groups}).
Zmenšení sektorových odchylek směrem k~AT nebo DK modelu je proto přímo cílem doporučené reformy: erga omnes s~reálnými minimálními tarify v~každém \ac{NACE} odvětví by stlačilo dolní odchylku bez nutnosti administrativního zasahování do horního pásma.

Graf odchylek sektorových mezd od celostátního mediánu (pro CZ vs.~AT a~DK) kvantifikuje míru polarizace trhu práce: v~ČR jsou odchylky výrazně větší --- ICT a~finance vykazují prémii +\SIrange{60}{80}{\percent} nad mediánem, zatímco obchod a~pohostinství zaostávají o~\SIrange{30}{40}{\percent}.
V~zemích s~erga omnes kolektivním vyjednáváním (AT, DK) je inter-sektorová disperze výrazně nižší: sektorové smlouvy s~celoodvětvovou platností zakotvují spodní sektorový standard, který redukuje rozsah, do jakého může spodní sektor klesat pod celostátní průměr.
Vysoká inter-sektorová disperze v~ČR má reálné sociální důsledky: pracovníci v~obchodu a~pohostinství --- skupiny s~vysokým podílem žen a~nízkokvalifikovaných --- jsou mzdově vytlačeni do in-work poverty territory (srov.~obr.~\ref{fig:arope_groups}).
Zmenšení sektorových odchylek směrem k~AT nebo DK modelu je proto přímo cílem doporučené reformy: erga omnes s~reálnými minimálními tarify v~každém \ac{NACE} odvětví by stlačilo dolní odchylku bez nutnosti administrativního zasahování do horního pásma.

Graf odchylek sektorových mezd od celostátního mediánu (pro CZ vs.~AT a~DK) kvantifikuje míru polarizace trhu práce: v~ČR jsou odchylky výrazně větší --- ICT a~finance vykazují prémii +\SIrange{60}{80}{\percent} nad mediánem, zatímco obchod a~pohostinství zaostávají o~\SIrange{30}{40}{\percent}.
V~zemích s~erga omnes kolektivním vyjednáváním (AT, DK) je inter-sektorová disperze výrazně nižší: sektorové smlouvy s~celoodvětvovou platností zakotvují spodní sektorový standard, který redukuje rozsah, do jakého může spodní sektor klesat pod celostátní průměr.
Vysoká inter-sektorová disperze v~ČR má reálné sociální důsledky: pracovníci v~obchodu a~pohostinství --- skupiny s~vysokým podílem žen a~nízkokvalifikovaných --- jsou mzdově vytlačeni do in-work poverty territory (srov.~obr.~\ref{fig:arope_groups}).
Zmenšení sektorových odchylek směrem k~AT nebo DK modelu je proto přímo cílem doporučené reformy: erga omnes s~reálnými minimálními tarify v~každém \ac{NACE} odvětví by stlačilo dolní odchylku bez nutnosti administrativního zasahování do horního pásma.

\input{texparts/python/sector_wages_deviation}



